\subsection{关于点测度积分的可测性}

命题 3. --- 设 \((\pi, g)\) 是一个 \(\mu\) -适应对，并令

\[
\nu = \int g (t) \varepsilon_ {\pi (t)} d \mu (t).
\]

设 \(f\) 是 \(X\) 到拓扑空间 \(G\) 中的一个映射，并设 \(S\) 是 \((\mu\text{-measurable})\) 集，由点 \(t \in T\) （满足 \(g(t) > 0\) 者）组成。\(f\) 为 \(\nu\) -可测的，必须且只须 \(f \circ \pi\) 在 \(S\) 上的限制是 \(\mu\) -可测的。

首先假设 \(f\) 是 \(\nu\) -可测的。由假设，集 \(\mathcal{K}\) （由紧子集 \(K\) （\(S\) 中的）组成，使 \(\pi\) 到 \(K\) 上的限制连续）是 \(\mu\) -稠密于 \(S\) 的（Ch. IV, §5, No.~10, 命题 15）。为证明 \(f \circ \pi\) 在 \(S\) 上的限制是 \(\mu\) -可测的，因此只需证明：对每个 \(K \in \mathcal{K}\) ，由紧子集 \(H\) （\(K\) 中的、使 \(f \circ \pi\) 在 \(H\) 上的限制连续者）组成的集是 \(\mu\) -稠密于 \(K\) 的（Ch. IV, §5, No.~8, 命题 13）。但由假设，存在紧集 \(\pi(K)\) 的一个分割，由一个 \(\nu\) -可忽略集 \(N\) 和一列紧集 \((C_n)\) 组成，使得 \(f\) 在每个 \(C_n\) 上的限制连续。在这些条件下，\(K \cap \overline{\pi}^{-1}(N)\) 与诸集 \(K \cap \overline{\pi}^{-1}(C_n)\) 构成 \(K\) 的一个分割；但 \(K \cap \overline{\pi}^{-1}(N)\) 依据 No.~2 的定理 1 的推论是 \(\mu\) -可忽略的，诸集 \(K \cap \overline{\pi}^{-1}(C_n)\) 是紧的，且 \(f \circ \pi\) 在后面这些集的每一个上的限制都连续，这就证明了 \(f \circ \pi\) 在 \(S\) 上的限制是 \(\mu\) -可测的。

反之，设情形如此；为证明 \(f\) 是 \(\nu\) -可测的，只需证明集 \(\mathfrak{L}\) （由紧子集 \(L\) （\(X\) 中的）组成，使 \(f\) 在 \(L\) 上的限制连续）是 \(\nu\) -稠密的（Ch. IV, §5, No.~10, 命题 15）。设 \(N\) 是 \(X\) 的一个子集，使 \(N \cap L\) 是 \(\nu\) -可忽略的（对每个 \(L \in \mathfrak{L}\) ），让我们证明 \(N\) 是局部 \(\nu\) -可忽略的。为此，我们必须证明 \(\overline{\pi}^{-1}(N) \cap S\) 是局部 \(\mu\) -可忽略的（No.~2 的定理 1 的推论）。现在，集 \(\mathfrak{H}\) （由紧子集 \(H\) （\(S\) 中的）组成，这些 \(H\) 使 \(\pi\) 和 \(f \circ \pi\) 在其上的限制都连续）由假设是 \(\mu\) -稠密于 \(S\) 的（Ch. IV, §5, No.~10, 命题 15）。因此只需证明 \(\overline{\pi}^{-1}(N) \cap H\) 是 \(\mu\) -可忽略的（对每个 \(H \in \mathfrak{H}\) ）。现在，\(\pi(H)\) 是紧的，并且可以与 \(H\) 按等价关系 \(\pi(t) = \pi(t')\) 所成的商空间等同，\(\pi\) 则等同于 \(H\) 到这个商空间上的典型映射（GT, I, §5, No.~2, 命题 3）。由于 \(f \circ \pi\) 在 \(H\) 上的限制连续，\(f\) 在 \(\pi(H)\) 上的限制就连续，换句话说 \(\pi(H) \in \mathfrak{L}\) ，从而 \(N \cap \pi(H)\) 是 \(\nu\) -可忽略的。由 No.~2 的定理 1 的推论，\(\overline{\pi}^{-1}(N \cap \pi(H)) \cap S\) 是局部 \(\mu\) -可忽略的；因此对集

\[
\mathrm{H} \cap \overline {{\pi}} ^ {- 1} (\mathrm{N}) \subset \overline {{\pi}} ^ {- 1} \bigl (\mathrm{N} \cap \pi (\mathrm{H}) \bigr) \cap \mathrm{S};
\]

但由于 H 是紧的，\(\mathrm{H} \cap \overline{\pi}^{-1}(\mathrm{N})\) 是 \(\mu\) -可忽略的，这就完成了证明。

注. --- 如果 f 是 X 到 Banach 空间 F 中的一个映射，那么说 \(f \circ \pi\) 在 S 上的限制是 \(\mu\) -可测的，与说函数 \((\mathbf{f} \circ \pi)g\) （定义在 T 上）是 \(\mu\) -可测的，是一回事，因为 \(g\) 是 \(\mu\) -可测的、在 S 中不为零且在 T - S 上为零（Ch. IV, §5, No.~10, 命题 15）。

